\section{Conclusion}
\label{sec:conclusion}

This project demonstrates a \gls{p2p} version of Battleship in which \glspl{zkp} replace the trusted game server for central fairness checks. Our scheme requires each player to commit to a hidden board state, prove that the initial board is valid, and then prove that every single reported $\hit$ or $\miss$ are consistent with the same commitment. The main contributions are the Circom/$\groth$ circuits for board and move verification, the \gls{p2p} proof\-/exchange structure, and the comparison between a direct $\poseidon$ commitment and a Merkle tree commitment. The Merkle tree approach is the scalable construction because it supports full\-/board commitments while requiring only a local authentication path for each move.

One direction for future work is to replace $\groth$ with a transparent proof system such as a \gls{stark}. A \gls{snark}, such as $\groth$, gives very short proofs and fast verification, but usually rely on elliptic\-/curve pairings and setup parameters tied to the circuit being proved. Instead, a \gls{stark} is built from hash functions and polynomial testing rather than from a trusted setup based on pairings. Thus, there is no hidden setup secret. Replacing $\groth$ with a \gls{stark} could address this limitation of the trusted setup directly, but at the potential cost of added latency.

Another direction is to redesign when and how move verification occurs. Instead of blocking the game on every proof, the protocol could allow the gameplay to continue optimistically while the proofs are generated and checked asynchronously. For instance, a move could be tentatively accepted, with the proof verified before the next state transition becomes final. If a proof fails or does not arrive before a timeout, the game could assign a forfeit. This change could hide some proving latency from the player and make the game work better in real\-/time scenarios.

A final direction is to strengthen the privacy guarantees for the attacking player. Garbled circuits \cite{yao86} and oblivious transfer \cite{rabin81,even85} are natural tools for this, because they allow two parties to compute a function over private inputs, the chosen shot and the committed board information, without revealing them directly. The output would reveal only the result $z$. Oblivious transfer would allow the attacker to obtain information corresponding to a chosen square without revealing which square was selected, while garbled circuits could enforce the computation $z$ itself. This would move the project from proving honest responses about public guesses toward a stronger protocol for hidden\-/information games where both the board and the targeting strategy remain private.
